%!TEX program = lualatex
\documentclass[a4paper,12pt]{article}

% =========================================================
% Idioma e tipografia - LuaLaTeX
% =========================================================
\usepackage[brazilian]{babel}
\usepackage{fontspec}
\setmainfont{Latin Modern Roman}

% =========================================================
% Layout
% =========================================================
\usepackage[top=2.4cm,bottom=2.3cm,left=2.5cm,right=2.5cm]{geometry}
\usepackage{setspace}
\setstretch{1.15}
\usepackage{indentfirst}
\setlength{\parindent}{1.0cm}
\setlength{\parskip}{0.25em}

% =========================================================
% Matemática, química e tabelas
% =========================================================
\usepackage{amsmath,amssymb}
\usepackage{siunitx}
\sisetup{output-decimal-marker={,},per-mode=symbol}
\usepackage[version=4]{mhchem}
\usepackage{booktabs}
\usepackage{tabularx}
\usepackage{array}
\usepackage{enumitem}

% =========================================================
% Caixas e cores
% =========================================================
\usepackage{xcolor}
\usepackage[most]{tcolorbox}
\definecolor{azulestudo}{RGB}{37,84,132}
\definecolor{laranjaestudo}{RGB}{202,112,27}
\definecolor{verdeestudo}{RGB}{52,112,73}
\definecolor{cinzaestudo}{RGB}{245,245,245}

\newtcolorbox{memorize}{
  breakable,
  colback=azulestudo!5,
  colframe=azulestudo,
  title=\textbf{Ponto para memorizar},
  fonttitle=\bfseries
}

\newtcolorbox{procedimento}{
  breakable,
  colback=verdeestudo!5,
  colframe=verdeestudo,
  title=\textbf{Procedimento},
  fonttitle=\bfseries
}

\newtcolorbox{atencao}{
  breakable,
  colback=laranjaestudo!7,
  colframe=laranjaestudo,
  title=\textbf{Atenção},
  fonttitle=\bfseries
}

\newtcolorbox{conta}{
  breakable,
  colback=cinzaestudo,
  colframe=black!55,
  title=\textbf{Conta essencial},
  fonttitle=\bfseries
}

% =========================================================
% Cabeçalho
% =========================================================
\usepackage{fancyhdr}
\pagestyle{fancy}
\fancyhf{}
\lhead{Química VIII -- Oceanografia}
\rhead{Guia de estudo -- prova de laboratório}
\cfoot{\thepage}
\setlength{\headheight}{15pt}

\usepackage[hidelinks]{hyperref}

\begin{document}

\begin{center}
    {\LARGE\bfseries Guia de Estudo -- Prova de Laboratório}\\[0.4em]
    {\large Química VIII -- Oceanografia}\\[0.4em]
    {\normalsize Conteúdo da apostila até a página 26}
\end{center}

\vspace{0.5em}

\begin{atencao}
Este material foi organizado para revisão da prova prática, com prioridade para os pontos explicitamente indicados: \textbf{segurança}, \textbf{procedimentos de pesagem}, \textbf{uso correto de material volumétrico}, \textbf{preparo de soluções}, \textbf{titulação/padronização} e \textbf{cálculos}. O conteúdo-base é o das páginas 3 a 26 da apostila de Química VIII.
\end{atencao}

\tableofcontents
\newpage

% =========================================================
\section{Segurança no laboratório}
% =========================================================

A segurança não deve ser estudada apenas como uma lista de proibições. Em prova prática, é comum a pergunta ser formulada como uma situação: \emph{``o que você faria antes de...?'', ``qual é a ordem correta?'', ``qual é o erro neste procedimento?''}. Portanto, é importante saber a regra e reconhecer sua aplicação.

\subsection{Regras que merecem atenção especial}

\begin{itemize}[leftmargin=1.6em]
    \item Seguir as instruções do professor ou técnico e \textbf{não trabalhar sozinho}.
    \item Comunicar imediatamente qualquer acidente.
    \item Não comer, beber, fumar ou brincar no laboratório.
    \item Usar guarda-pó comprido, fechado e de algodão, além de \textbf{óculos de segurança}.
    \item Não usar roupas ou calçados que deixem as pernas e pés expostos; cabelos longos devem permanecer presos.
    \item Não inalar nem provar reagentes.
    \item Conferir o \textbf{rótulo} antes de usar qualquer reagente e manter os frascos fechados quando não estiverem em uso.
    \item Manter solventes inflamáveis afastados de chama, forno, aquecimento e luz solar direta.
    \item Não aquecer líquidos inflamáveis em chama direta.
    \item Vidro quente tem aparência semelhante à de vidro frio.
    \item Nunca pipetar líquidos cáusticos ou tóxicos com a boca: usar pipetador ou pera de borracha.
    \item Não descartar sólidos na pia. O descarte de reagentes e materiais deve seguir a orientação do técnico.
    \item Ao terminar, lavar e guardar a vidraria, limpar a bancada e verificar torneiras e gás.
\end{itemize}

\begin{memorize}
\textbf{ÁCIDO NA ÁGUA, NUNCA ÁGUA NO ÁCIDO.}\\
Ao preparar uma solução aquosa de ácido, primeiro deve existir água no recipiente; o ácido concentrado é então adicionado \textbf{lentamente} à água.
\end{memorize}

\begin{procedimento}
\textbf{Exemplo da própria prática de HCl:}
\begin{enumerate}[leftmargin=1.7em]
    \item colocar cerca de \SI{250}{\milli\liter} de água destilada no recipiente;
    \item adicionar lentamente o HCl concentrado;
    \item esperar o sistema resfriar;
    \item somente depois completar o volume com água destilada e homogeneizar.
\end{enumerate}
\end{procedimento}

\paragraph{Explicação complementar útil para a prova.}
A diluição de um ácido concentrado libera calor. Colocar o ácido lentamente em uma quantidade maior de água favorece a dissipação desse calor e reduz o risco de aquecimento localizado e respingos. Esta explicação é complementar; a regra operacional da apostila é: \textbf{ácido concentrado na água, nunca o contrário}.

% =========================================================
\section{Técnicas básicas de laboratório}
% =========================================================

\subsection{Seleção de reagentes}

A pureza do reagente deve ser compatível com o uso analítico. A apostila destaca:

\begin{itemize}[leftmargin=1.6em]
    \item \textbf{P.A. (pró-análise):} reagentes que atendem requisitos mínimos de pureza e são usados, sempre que possível, em trabalhos analíticos;
    \item \textbf{reagentes de uso específico:} preparados para técnicas determinadas, como grau HPLC ou padrões para espectrometria atômica.
\end{itemize}

\subsection{Limpeza e identificação}

Como as análises são frequentemente realizadas em duplicata ou triplicata, os frascos precisam ser identificados. A sequência de limpeza apresentada é:

\begin{procedimento}
\begin{enumerate}[leftmargin=1.7em]
    \item lavar com detergente;
    \item enxaguar exaustivamente com água corrente;
    \item fazer a \textbf{rinsagem final} com água destilada e/ou deionizada.
\end{enumerate}
Em uma vidraria perfeitamente limpa, gotas de água não devem permanecer aderidas às paredes.
\end{procedimento}

\subsection{Evaporação}

A evaporação pode ser usada para reduzir o volume de uma solução contendo soluto não volátil. O arranjo mostrado na apostila usa um \textbf{vidro de relógio apoiado em bastão de vidro}: o vapor consegue sair, enquanto o sistema fica parcialmente protegido contra contaminação e respingos.

\subsection{Dessecador e dessecantes}

Materiais secos em estufa podem reabsorver umidade do ar. Por isso, são mantidos em dessecador. A apostila cita como dessecantes \ce{CaCl2} anidro, \ce{CaSO4}, \ce{Mg(ClO4)2} e sílica-gel.

\subsection{Filtração e lavagem}

A filtração é uma separação sólido--líquido, importante quando se forma um precipitado.

\begin{memorize}
A sequência destacada é: \textbf{decantação $\rightarrow$ lavagem $\rightarrow$ transferência do precipitado}.
\end{memorize}

Cuidados relevantes:

\begin{itemize}[leftmargin=1.6em]
    \item usar bastão de vidro para dirigir o fluxo;
    \item encostar a extremidade da haste do funil na parede interna do béquer que recebe o filtrado;
    \item não deixar o precipitado secar no filtro durante a filtração;
    \item manter o nível da solução em aproximadamente $3/4$ da altura do papel de filtro;
    \item umedecer o papel de filtro para que ele fique bem aderido ao funil;
    \item escolher papel de filtro cujo topo fique cerca de \SIrange{1}{2}{\centi\meter} abaixo da borda do funil.
\end{itemize}

\subsection{Secagem e calcinação}

\begin{center}
\begin{tabularx}{\textwidth}{>{\bfseries}p{3.2cm} X X}
\toprule
 & Secagem & Calcinação \\
\midrule
Temperatura & inferior a \SI{250}{\degreeCelsius} & superior a \SI{250}{\degreeCelsius} \\
Objetivo típico & remover água residual de lavagem & promover transformação/decomposição térmica até a forma final a ser pesada \\
Equipamento & estufa & mufla ou aquecimento direto, conforme o procedimento \\
\bottomrule
\end{tabularx}
\end{center}

% =========================================================
\section{Medidas de massa e balança analítica}
% =========================================================

A balança analítica é empregada quando se necessita de medida precisa de massa. A apostila chama atenção para \textbf{sensibilidade} e \textbf{exatidão}, além do controle de tara nas balanças eletrônicas.

\subsection{Principais fontes de erro}

\begin{center}
\begin{tabularx}{\textwidth}{>{\bfseries}p{3.5cm} X}
\toprule
Fonte de erro & Consequência / prevenção \\
\midrule
Carga estática & Vidro ou porcelana atritados com papel/tecido podem eletrizar-se. Deve-se aguardar dissipação da carga antes da pesagem. \\
Atmosfera & Alguns materiais ganham ou perdem massa por absorção/perda de água, \ce{CO2}, etc. \\
Objeto quente & Pode produzir correntes de convecção e alterar a leitura. Deve estar na mesma temperatura da balança. \\
Reabsorção de umidade & Objetos aquecidos devem resfriar em dessecador antes da medida. \\
Empuxo do ar & Produz massa aparente; em sólidos, a correção é geralmente pequena, mas pode ser importante para materiais de baixa densidade. \\
\bottomrule
\end{tabularx}
\end{center}

\subsection{Cuidados com a balança}

\begin{itemize}[leftmargin=1.6em]
    \item Colocar o objeto próximo ao \textbf{centro do prato}.
    \item Não tocar com as mãos o objeto a ser pesado; usar pinça ou papel limpo.
    \item O objeto deve estar à temperatura da balança.
    \item Nunca colocar reagente diretamente sobre o prato; usar vidro de relógio, pesa-filtro, béquer pequeno etc.
    \item Se houver derramamento, limpar imediatamente com pincel macio.
    \item Manter as portas da balança fechadas durante a leitura.
    \item Ao terminar, retirar o material e deixar a balança limpa e na condição indicada pelo equipamento.
\end{itemize}

\subsection{Métodos de pesagem}

\begin{center}
\begin{tabularx}{\textwidth}{>{\bfseries}p{3.7cm} X}
\toprule
Método & Uso indicado na apostila \\
\midrule
Pesagem direta & Determinar a massa de objetos ou materiais delimitados, como cadinho, vidro de relógio e pesa-filtro. \\
Pesagem por adição & Obter uma massa previamente estabelecida de sólido a granel ou líquido. \\
Pesagem por diferença & Usada para substâncias higroscópicas ou que se alterem em contato com o ar. \\
\bottomrule
\end{tabularx}
\end{center}

\subsection{Prática 1: pesagem por adição de NaCl}

\begin{procedimento}
Para obter aproximadamente \SI{0.5000}{\gram} de \ce{NaCl}:
\begin{enumerate}[leftmargin=1.7em]
    \item verificar se a balança está nivelada;
    \item ligar e verificar o zero;
    \item pesar diretamente um vidro de relógio limpo, seco e à temperatura ambiente: $m_1$;
    \item adicionar cuidadosamente o NaCl com espátula;
    \item fechar a porta da balança e registrar a massa do conjunto: $m_2$;
    \item retirar o conjunto e verificar se a balança ficou limpa.
\end{enumerate}
\end{procedimento}

\begin{conta}
A massa de NaCl é obtida por diferença:
\[
\boxed{m_{\ce{NaCl}} = m_2 - m_1}
\]
Embora o procedimento seja chamado de \textbf{pesagem por adição}, a massa efetivamente adicionada é determinada pela diferença entre a massa final do conjunto e a massa inicial do recipiente.
\end{conta}

% =========================================================
\section{Medidas de volume}
% =========================================================

A precisão das medidas volumétricas é tão importante quanto a precisão da massa. Os materiais volumétricos são calibrados a uma temperatura de referência, geralmente \SI{20}{\degreeCelsius}.

\subsection{Pipeta volumétrica}

\begin{procedimento}
Sequência de uso:
\begin{enumerate}[leftmargin=1.7em]
    \item limpar: detergente $\rightarrow$ água corrente $\rightarrow$ água destilada/deionizada;
    \item rinsar internamente com pequena quantidade da solução que será transferida;
    \item aspirar com pera de borracha até um pouco acima do traço;
    \item verificar se existem bolhas;
    \item ajustar lentamente o menisco ao traço de referência;
    \item transferir mantendo a pipeta na vertical;
    \item deixar a ponta encostada na parede do recipiente por cerca de \SI{15}{\second} após o escoamento.
\end{enumerate}
\end{procedimento}

\begin{memorize}
Na pipeta volumétrica comum, \textbf{não se sopra automaticamente a gota residual da ponta}. A apostila ressalta que somente alguns tipos especiais exigem a eliminação desse líquido residual.
\end{memorize}

\subsection{Bureta}

A bureta transfere volumes variáveis e é o equipamento fundamental para titulações.

\begin{procedimento}
Antes da titulação:
\begin{enumerate}[leftmargin=1.7em]
    \item garantir que a bureta esteja perfeitamente limpa;
    \item conferir se a torneira está fechada;
    \item colocar \SIrange{5}{10}{\milli\liter} do titulante, girar a bureta para molhar o interior e deixar essa porção escoar;
    \item preencher acima da marca de \SI{0.00}{\milli\liter};
    \item escoar até ajustar o nível inicial.
\end{enumerate}
Durante a titulação: manter o erlenmeyer em agitação; próximo ao ponto final, rinsar suas paredes com água destilada e adicionar o titulante \textbf{gota a gota}.
\end{procedimento}

\begin{memorize}
A leitura do nível da bureta deve ser feita com os \textbf{olhos no nível do menisco}, evitando erro de paralaxe. A apostila pede leitura do volume gasto com precisão de \SI{0.05}{\milli\liter} nos procedimentos de titulação descritos.
\end{memorize}

\subsection{Balão volumétrico}

O balão volumétrico é calibrado para \textbf{conter} um volume fixo quando o menisco está no traço.

\begin{procedimento}
Preparo de solução a partir de um soluto sólido:
\begin{enumerate}[leftmargin=1.7em]
    \item pesar o soluto em pequeno béquer ou pesa-filtro;
    \item dissolver usando o menor volume possível de solvente;
    \item transferir quantitativamente para o balão com funil/bastão;
    \item rinsar béquer, bastão e funil e transferir também as águas de lavagem;
    \item completar até o traço com água destilada;
    \item tampar e homogeneizar.
\end{enumerate}
\end{procedimento}

% =========================================================
\section{Prática 2: calibração de uma pipeta volumétrica}
% =========================================================

A calibração é feita medindo a massa de água entregue pela pipeta e convertendo essa massa em volume usando a densidade da água na temperatura medida.

Para a pipeta de \SI{25.00}{\milli\liter}, a apostila informa desvio máximo de
$\pm\SI{0.02}{\milli\liter}$. O tempo de escoamento indicado é de aproximadamente \SI{25}{\second}.

\begin{procedimento}
\begin{enumerate}[leftmargin=1.7em]
    \item medir a massa do erlenmeyer vazio: $m_1$;
    \item transferir a água da pipeta para o erlenmeyer;
    \item medir a massa do conjunto: $m_2$;
    \item calcular a massa de água;
    \item medir a temperatura da água e obter $\rho$ na tabela;
    \item calcular o volume da pipeta;
    \item repetir a determinação e usar a média.
\end{enumerate}
\end{procedimento}

\begin{conta}
\[
\boxed{m_{\ce{H2O}} = m_2-m_1}
\]
\[
\boxed{V = \frac{m_{\ce{H2O}}}{\rho_{\ce{H2O}}}}
\]
Se forem feitas duas determinações:
\[
\boxed{\overline{V}=\frac{V_1+V_2}{2}}
\]
Unidades: se $m$ estiver em g e $\rho$ em g/mL, $V$ sai em mL.
\end{conta}

% =========================================================
\section{Análise volumétrica e titulação}
% =========================================================

Na análise volumétrica, mede-se com exatidão o volume de uma solução de concentração conhecida necessário para reagir completamente com o analito.

\subsection{Classificação}

\begin{itemize}[leftmargin=1.6em]
    \item volumetria de neutralização;
    \item volumetria de precipitação;
    \item volumetria de complexação;
    \item volumetria de oxirredução.
\end{itemize}

\subsection{Requisitos da reação}

A reação usada deve ser rápida, completa, descrita por uma única equação ajustada, possuir estequiometria conhecida e permitir a determinação satisfatória do ponto final.

\begin{memorize}
Em uma titulação, não basta ``mudar de cor''. A mudança observada com o indicador é o sinal experimental usado para reconhecer o ponto final; o cálculo depende da \textbf{estequiometria da reação} e do \textbf{volume efetivamente gasto} do titulante.
\end{memorize}

% =========================================================
\section{Prática 3: preparo e padronização de HCl}
% =========================================================

A solução concentrada de HCl indicada na apostila tem aproximadamente \SI{37}{\percent} (m/m) e densidade aproximada de \SI{1.19}{\gram\per\milli\liter}. Ela \textbf{não é padrão primário}; por isso, prepara-se uma concentração aproximada por diluição e depois se faz a padronização.

\subsection{Cálculo geral do volume de HCl concentrado}

Para preparar uma solução de concentração $C_f$ e volume final $V_f$:

\begin{conta}
\textbf{1. Mols necessários:}
\[
\boxed{n=C_fV_f}
\]

\textbf{2. Massa de HCl puro:}
\[
\boxed{m_{\ce{HCl}}=n\,M_{\ce{HCl}}}
\]

\textbf{3. Massa de solução concentrada}, usando a fração mássica $w=T/100$:
\[
\boxed{m_{\mathrm{sol.conc.}}=\frac{m_{\ce{HCl}}}{w}}
\]

\textbf{4. Volume da solução concentrada:}
\[
\boxed{V_{\mathrm{sol.conc.}}=\frac{m_{\mathrm{sol.conc.}}}{\rho}}
\]
\end{conta}

\subsection{Exemplo: \SI{500}{\milli\liter} de HCl \SI{0.1}{\mole\per\liter}}

Usando $M_{\ce{HCl}}=\SI{36.5}{\gram\per\mole}$, $w=0{,}37$ e $\rho=\SI{1.19}{\gram\per\milli\liter}$:

\[
n=(0{,}1)(0{,}500)=0{,}050\ \mathrm{mol}
\]
\[
m_{\ce{HCl}}=(0{,}050)(36{,}5)=1{,}825\ \mathrm{g}
\]
\[
m_{\mathrm{sol.conc.}}=\frac{1{,}825}{0{,}37}\approx4{,}93\ \mathrm{g}
\]
\[
V_{\mathrm{sol.conc.}}=\frac{4{,}93}{1{,}19}\approx\boxed{4{,}14\ \mathrm{mL}}
\]

\begin{atencao}
\textbf{Inconsistência numérica da apostila:} o texto da Prática 3 declara o preparo de \SI{500}{\milli\liter} de HCl \SI{0.1}{\mole\per\liter}, mas o cálculo intermediário impresso usa \SI{1000}{\milli\liter} e chega a \SI{3.65}{\gram} de HCl puro. Para \SI{500}{\milli\liter}, o cálculo estequiométrico direto fornece \SI{1.825}{\gram} de HCl puro e, com os valores aproximados de 37\% e \SI{1.19}{\gram\per\milli\liter}, cerca de \SI{4.14}{\milli\liter} do ácido concentrado. Em prova, confira se o enunciado pede \SI{500}{\milli\liter} ou \SI{1}{\liter} antes de substituir os números.
\end{atencao}

\subsection{Procedimento de preparo do HCl}

\begin{procedimento}
\begin{enumerate}[leftmargin=1.7em]
    \item verificar no rótulo a densidade e a percentagem mássica do HCl concentrado;
    \item calcular o volume necessário;
    \item medir o ácido usando o material indicado e fazer a transferência \textbf{na capela}, com exaustão ligada e pipetador/bulbo;
    \item colocar cerca de \SI{250}{\milli\liter} de água destilada no recipiente de \SI{500}{\milli\liter};
    \item adicionar o ácido concentrado \textbf{lentamente à água};
    \item esperar resfriar;
    \item completar o volume com água destilada e homogeneizar;
    \item transferir para frasco limpo e previamente rinsado;
    \item identificar com nome da solução, concentração, fator, data e analista.
\end{enumerate}
\end{procedimento}

\subsection{Padronização do HCl com bórax}

O bórax, \ce{Na2B4O7.10H2O}, é usado como padrão primário alcalino. A apostila fornece massa molar \SI{381.37}{\gram\per\mole}.

A reação global apresentada é:
\[
\ce{Na2B4O7 + 2HCl + 5H2O -> 2NaCl + 4H3BO3}
\]
Portanto, \textbf{1 mol de bórax corresponde a 2 mol de HCl}.

\begin{conta}
Concentração do bórax preparado a partir de massa $m$ e volume $V$ em litros:
\[
\boxed{C_{\mathrm{borax}}=\frac{m}{381{,}37\,V}}
\]

Correção pelo ensaio em branco:
\[
\boxed{V_{\ce{HCl}}=V_t-V_B}
\]

Concentração do HCl:
\[
\boxed{C_{\ce{HCl}}=\frac{2\,C_{\mathrm{borax}}\,V_{\mathrm{borax}}}{V_{\ce{HCl}}}}
\]
Os volumes do numerador e denominador podem estar ambos em mL, pois a razão cancela a unidade de volume.

Resultado final de duas determinações:
\[
\boxed{\overline{C}_{\ce{HCl}}=\frac{C_{\ce{HCl},1}+C_{\ce{HCl},2}}{2}}
\]
\end{conta}

\begin{procedimento}
Titulação pelo bórax:
\begin{enumerate}[leftmargin=1.7em]
    \item transferir \SI{50.00}{\milli\liter} da solução de bórax para o erlenmeyer;
    \item adicionar 2--3 gotas de alaranjado de metila;
    \item rinsar a bureta com o HCl preparado;
    \item preencher a bureta com HCl;
    \item colocar o erlenmeyer sobre fundo branco;
    \item adicionar ácido lentamente, agitando;
    \item perto da virada, lavar as paredes do erlenmeyer com água destilada e prosseguir gota a gota;
    \item registrar $V_t$;
    \item corrigir pelo branco: $V_{\ce{HCl}}=V_t-V_B$;
    \item repetir com outra alíquota.
\end{enumerate}
\end{procedimento}

% =========================================================
\section{Prática 4: preparo e padronização de NaOH}
% =========================================================

O NaOH não é padrão primário. A apostila destaca dois problemas: \textbf{pureza do reagente} e \textbf{reação com o \ce{CO2} do ar}.

\subsection{Efeito do \texorpdfstring{\ce{CO2}}{CO2}}

\[
\ce{CO2(g) + 2OH^- -> CO3^{2-} + H2O}
\]

A presença de carbonato altera o comportamento da titulação dependendo do indicador. A apostila explica que, com indicador na faixa alcalina, como fenolftaleína, o carbonato reage com apenas um \ce{H3O+} até o ponto observado, podendo diminuir a concentração efetiva da base e produzir erro sistemático.

\begin{memorize}
Para reduzir a interferência de \ce{CO2}:
\begin{itemize}[leftmargin=1.6em]
    \item a água usada no preparo deve ser pobre em \ce{CO2};
    \item se for usada água destilada, a apostila recomenda fervê-la para remover o gás e depois deixá-la \textbf{resfriar antes de adicionar a base};
    \item a água deionizada é preferida no procedimento descrito;
    \item a solução padronizada não deve permanecer longos períodos em frasco de vidro; a apostila recomenda frasco de polietileno para armazenamento da solução preparada.
\end{itemize}
\end{memorize}

\subsection{Preparo da solução de NaOH}

\begin{procedimento}
\begin{enumerate}[leftmargin=1.7em]
    \item pesar \SI{50}{\gram} de NaOH em béquer de \SI{250}{\milli\liter};
    \item dissolver em \SI{50}{\milli\liter} de água destilada isenta de \ce{CO2};
    \item transferir para proveta de \SI{100}{\milli\liter} e deixar o precipitado sedimentar por \SI{24}{\hour};
    \item retirar cuidadosamente \SI{3.20}{\milli\liter} da solução concentrada e límpida;
    \item transferir para recipiente de \SI{500}{\milli\liter} que já contenha cerca de \SI{250}{\milli\liter} de água destilada;
    \item aguardar resfriamento e completar o volume com água isenta de \ce{CO2};
    \item homogeneizar;
    \item transferir para frasco de polietileno previamente lavado e rinsado com a própria solução;
    \item identificar com nome, concentração molar, fator, data e analista.
\end{enumerate}
\end{procedimento}

\subsection{Padronização do NaOH com HCl}

A reação é 1:1:
\[
\ce{HCl + NaOH -> NaCl + H2O}
\]

\begin{conta}
Após corrigir o volume de HCl pelo ensaio em branco:
\[
V_{\ce{HCl}}=V_t-V_B
\]

A concentração do NaOH é:
\[
\boxed{C_{\ce{NaOH}}=\frac{C_{\ce{HCl}}\,V_{\ce{HCl}}}{V_{\ce{NaOH}}}}
\]

Para duas determinações:
\[
\boxed{\overline{C}_{\ce{NaOH}}=\frac{C_{\ce{NaOH},1}+C_{\ce{NaOH},2}}{2}}
\]
\end{conta}

\begin{procedimento}
\begin{enumerate}[leftmargin=1.7em]
    \item transferir \SI{25}{\milli\liter} de NaOH para erlenmeyer;
    \item adicionar \SI{25}{\milli\liter} de água destilada e 3--4 gotas de alaranjado de metila;
    \item rinsar a bureta com HCl padrão e preenchê-la;
    \item colocar o erlenmeyer sobre fundo branco;
    \item titular lentamente, agitando;
    \item próximo à virada, lavar as paredes do erlenmeyer e adicionar HCl gota a gota;
    \item registrar o volume gasto, corrigir pelo branco e repetir com nova alíquota.
\end{enumerate}
\end{procedimento}

\begin{atencao}
Na página 26, o item 5.2 contém a frase ``anotar o volume de base consumido'', mas o próprio procedimento estabelece que a \textbf{bureta contém HCl} e, no item 5.8, manda registrar o \textbf{volume de ácido consumido}. Para estudar o procedimento, considere o HCl como titulante e o NaOH como solução no erlenmeyer.
\end{atencao}

% =========================================================
\section{Roteiros de cálculo para a prova}
% =========================================================

\subsection{Tipo A: ``prepare uma solução a partir de reagente sólido''}

Se o soluto sólido tiver pureza tratada como adequada para o cálculo:
\[
n=CV
\]
\[
m=nM
\]
Logo:
\[
\boxed{m=CVM}
\]
Depois: pesar $\rightarrow$ dissolver $\rightarrow$ transferir quantitativamente $\rightarrow$ completar até o traço $\rightarrow$ homogeneizar.

\subsection{Tipo B: ``prepare uma solução por diluição de reagente concentrado''}

Se forem dados densidade e teor em massa:
\[
\boxed{n=C_fV_f}
\]
\[
\boxed{m_{\mathrm{soluto}}=nM}
\]
\[
\boxed{m_{\mathrm{solucao}}=\frac{m_{\mathrm{soluto}}}{T/100}}
\]
\[
\boxed{V_{\mathrm{concentrado}}=\frac{m_{\mathrm{solucao}}}{\rho}}
\]

\subsection{Tipo C: ``calibre a pipeta''}

\[
\boxed{m_{\ce{H2O}}=m_2-m_1}
\qquad
\boxed{V=\frac{m}{\rho}}
\]

\subsection{Tipo D: ``padronize o HCl com bórax''}

\[
\boxed{C_{\mathrm{borax}}=\frac{m}{381{,}37V}}
\]
\[
\boxed{V_{\ce{HCl}}=V_t-V_B}
\]
\[
\boxed{C_{\ce{HCl}}=\frac{2C_{\mathrm{borax}}V_{\mathrm{borax}}}{V_{\ce{HCl}}}}
\]

\subsection{Tipo E: ``padronize o NaOH com HCl''}

\[
\boxed{V_{\ce{HCl}}=V_t-V_B}
\]
\[
\boxed{C_{\ce{NaOH}}=\frac{C_{\ce{HCl}}V_{\ce{HCl}}}{V_{\ce{NaOH}}}}
\]

% =========================================================
\section{Perguntas de revisão rápida}
% =========================================================

\begin{enumerate}[leftmargin=1.8em]
    \item \textbf{Ao preparar ácido diluído, o que entra primeiro no recipiente?}\\
    Água. O ácido concentrado é adicionado lentamente à água.

    \item \textbf{Por que se espera o HCl diluído resfriar antes de completar o volume?}\\
    O procedimento da apostila exige o resfriamento antes da aferição final; como explicação complementar, a diluição libera calor e a temperatura afeta o volume.

    \item \textbf{Pode pipetar HCl com a boca?}\\
    Não. Deve-se usar pipetador ou pera de borracha, e o procedimento com HCl concentrado é realizado em capela.

    \item \textbf{Qual é a sequência de limpeza da vidraria?}\\
    Detergente $\rightarrow$ água corrente $\rightarrow$ água destilada/deionizada.

    \item \textbf{O que indica uma vidraria perfeitamente limpa?}\\
    Gotas de água não ficam retidas nas paredes.

    \item \textbf{Qual a diferença entre pesagem direta, por adição e por diferença?}\\
    Direta: massa do objeto/material delimitado. Adição: obtenção de massa previamente estabelecida. Diferença: indicada para substâncias higroscópicas ou alteráveis pelo ar.

    \item \textbf{Por que não se pesa um objeto quente?}\\
    Porque a diferença de temperatura pode gerar correntes de convecção e erro de leitura.

    \item \textbf{Por que não se toca com a mão no recipiente que será pesado?}\\
    Para evitar contaminação/engorduramento e alteração da massa medida.

    \item \textbf{Como se lê corretamente uma bureta?}\\
    Com os olhos no nível do menisco.

    \item \textbf{O que se faz perto do ponto final de uma titulação?}\\
    Rinsa-se a parede do erlenmeyer com água destilada, mantém-se a agitação e adiciona-se o titulante gota a gota.

    \item \textbf{Por que a bureta é rinsada com o titulante antes do enchimento?}\\
    O procedimento da apostila exige a rinsagem com a própria solução que será usada na titulação, evitando que água residual permaneça como meio de diluição do titulante.

    \item \textbf{Para que serve o ensaio em branco?}\\
    Para corrigir o volume medido: $V_{\mathrm{real}}=V_t-V_B$.

    \item \textbf{Por que o HCl preparado precisa ser padronizado?}\\
    Porque o HCl concentrado usado na diluição não é padrão primário; a solução preparada é inicialmente apenas aproximada.

    \item \textbf{Por que o NaOH precisa ser padronizado?}\\
    Porque não possui pureza de padrão primário e pode reagir com \ce{CO2} atmosférico.

    \item \textbf{Qual é a estequiometria entre bórax e HCl na reação global?}\\
    1 mol de bórax para 2 mol de HCl.

    \item \textbf{Qual é a estequiometria entre HCl e NaOH?}\\
    1:1.
\end{enumerate}

% =========================================================
\section{Checklist final para a véspera}
% =========================================================

\begin{memorize}
Antes da prova, verifique se você consegue explicar sem consultar:
\begin{itemize}[leftmargin=1.6em]
    \item ácido na água, nunca o contrário;
    \item sequência de limpeza e rinsagem;
    \item pesagem direta, por adição e por diferença;
    \item cuidados com a balança e principais erros de massa;
    \item uso correto de pipeta, bureta e balão volumétrico;
    \item leitura do menisco sem paralaxe;
    \item calibração da pipeta por $V=m/\rho$;
    \item cálculo de solução a partir de $n=CV$;
    \item cálculo de HCl concentrado usando teor mássico e densidade;
    \item preparo seguro do HCl;
    \item por que HCl e NaOH são padronizados;
    \item correção por ensaio em branco;
    \item fator estequiométrico 2 entre bórax e HCl;
    \item relação 1:1 entre HCl e NaOH;
    \item efeito de \ce{CO2} sobre soluções de NaOH;
    \item armazenamento do NaOH em frasco de polietileno no procedimento da apostila.
\end{itemize}
\end{memorize}

\vfill
\begin{center}
\small\textit{Material de estudo elaborado a partir da apostila de aulas práticas de Química VIII -- Oceanografia, conteúdo até a página 26.}
\end{center}

\end{document}
